\section{Methodology}
\label{sec:methodology}

The objective of this research is to establish a rigorous, high-precision conversion framework between barometric pressure and geometric altitude in complex urban environments. Unlike pure black-box deep learning approaches, our methodology decomposes the altitude measurement into a deterministic physical baseline and a learnable spatial residual component. To achieve this, we co-design an edge-sensing hardware terminal and a Physics-Informed Neural Field (PINF) algorithm.

\begin{figure}[htbp]
    \centering
    % Name your image fig2_hardware.png and put it in the image folder
    \includegraphics[width=\columnwidth]{image/figure2.pdf}
    \caption{Hardware architecture and signal flow of the GeoBox data acquisition terminal. The system integrates GNSS, barometric, and environmental sensors. An onboard MCU performs strict temporal synchronization and simple sensor fusion before transmitting the data payload via a 4G LTE/MQTT interface.}
    \label{fig:hardware}
\end{figure}

\subsection{Edge Sensing Instrumentation}
\label{subsec:hardware}

The foundation of our proposed framework is a custom-designed, low-cost multi-sensor data acquisition terminal, designated as the ``GeoBox''. This edge device is engineered to capture synchronized navigation and atmospheric profiles across distributed urban topologies, providing the sparse but high-fidelity data required to train the continuous neural field. The hardware architecture integrates three core sensing modules:
\begin{enumerate}
    \item \textbf{GNSS Positioning Module}: An integrated GNSS receiver coupled with an external active antenna acquires high-precision geometric coordinates (latitude $\phi$, longitude $\lambda$) and the true Geometric Altitude ($h_\text{GNSS}$) relative to the WGS-84 reference ellipsoid. It supports multiple constellations (GPS, BeiDou, Galileo) to ensure lock stability in urban canyons.
    \item \textbf{Barometric Pressure Sensor}: A high-resolution MEMS-based static pressure sensor provides raw atmospheric pressure measurements ($P_\text{obs}$) with an operating range of 300 to 1100 hPa. The sensor is isolated within a specially designed enclosure featuring vented static ports to minimize dynamic pressure artifacts induced by wind gusts or platform motion.
    \item \textbf{Environmental Sensing Array}: A digital hygrometer and thermometer are integrated to measure ambient temperature ($T_\text{obs}$) and relative humidity ($RH$). These parameters are strictly required to calculate the specific humidity and virtual temperature, effectively compensating for air density variations in the subsequent hydrostatic calculations.
\end{enumerate}

\subsubsection{Data Synchronization and Transmission}
Metrological integrity depends heavily on the temporal synchronization of the heterogeneous data streams. An onboard Microcontroller Unit (MCU) performs time synchronization, ensuring that the GNSS geometric altitude is strictly time-stamped alongside the barometric and environmental readings. To guarantee data reliability during field deployments, the system executes a Built-In Test (BIT) sequence upon startup to verify GNSS lock status, network connectivity, and sensor health. Powered by a lithium-ion battery for extended autonomy, the MCU aggregates the data into a standardized payload—comprising a unique device ID (UID), Unix timestamp, $P_\text{obs}$, $T_\text{obs}$, $RH$, and $(\phi, \lambda)$ coordinates. This packet is transmitted in real-time to a centralized cloud server via a 4G/LTE cellular module utilizing the MQTT protocol. This continuous, synchronized data pipeline forms the empirical basis for the physical baseline estimation and neural residual learning phases.

\subsection{Physical Baseline Estimation and Problem Formulation}
\label{subsec:physical_baseline}

To isolate the complex, non-linear microclimate disturbances inherent to urban VLL airspace, we first establish a deterministic physical baseline. This baseline serves as the mean function for our estimation framework, derived from fundamental atmospheric physics and the edge-sensor measurements acquired by the GeoBox.

The foundational relationship between atmospheric pressure $P$ and geopotential height $H$ is governed by the hydrostatic equation:
\begin{equation}
\frac{dP}{dH} = -\rho g \ ,
\label{eq:hydrostatic}
\end{equation}
where $\rho$ is the air density and $g$ is the gravitational acceleration. Standard aviation models (e.g., ISA) assume a static dry-air atmosphere \cite{icao1993manual}. However, to achieve meter-level accuracy, variations in air density caused by moisture must be accounted for. We utilize the onboard environmental sensors to compute the in-situ Virtual Temperature $T_\text{v}$:
\begin{equation}
T_\text{v} = T_\text{obs} (1 + 0.61 q) \ ,
\label{eq:virtual_temp}
\end{equation}
where $T_\text{obs}$ is the measured absolute temperature and $q$ is the specific humidity derived from the measured relative humidity ($RH$) \cite{simonetti2024geodetic}.

By substituting the ideal gas law ($\rho = P / (R_\text{dry} T_\text{v})$) into Eq. (\ref{eq:hydrostatic}) and integrating, we obtain the Hypsometric Equation, which defines our physical baseline estimator $\hat{H}_\text{phy}$:
\begin{equation}
\hat{H}_\text{phy} = H_\text{ref} + \frac{R_\text{dry} T_\text{v}}{g_0} \ln\left(\frac{P_\text{ref}}{P_\text{obs}}\right) \ ,
\label{eq:hypsometric}
\end{equation}
where $R_\text{dry}$ is the specific gas constant for dry air, $g_0$ is the standard gravity, and $P_\text{ref}$ and $H_\text{ref}$ are the reference pressure and height, respectively. This step accounts for the primary altitude variance driven by macro-meteorological conditions.

Crucially, $\hat{H}_\text{phy}$ represents the orthometric height (relative to the global mean sea level, MSL). To harmonize this standard atmospheric altitude with standard UAS GNSS coordinate frames—which typically output Height Above Ellipsoid (HAE)—one must traditionally apply the Geoid Undulation anomaly $N(\phi, \lambda)$ retrieved from high-resolution gravity models (e.g., EGM96/EGM2008).
However, deploying global high-resolution geoid models heavily limits real-time edge processing capabilities. In our proposed GeoBox edge-sensing architecture, the integrated multi-constellation GNSS receiver is uniquely engineered to output true orthometric heights ($h_\text{GNSS} \equiv$ MSL) directly by utilizing onboard regional geoid anomaly grids. Therefore, we utilize the mathematically rigorous $\hat{H}_\text{phy}$ directly as the un-deformed physical baseline $h_\text{phy}$, decoupling the high-frequency measurement calibration from expensive offline gravitational geoid interpolations.

While $h_\text{phy}$ provides a strong generalized mean function, it fundamentally fails to capture high-frequency local atmospheric disturbances, such as urban heat islands, building wake turbulence, and importantly, the individual static bias inherent to uncalibrated hardware, which cannot be modeled by simple analytical equations. To enforce physical consistency across multiple uncalibrated nodes, we introduce a Physics-Derived Pressure Bias Formulation. Instead of directly predicting an explicit geometric residual ($\Delta h$), we predict the localized barometric pressure correction ($\delta P$).

The unified urban altitude estimation problem is thus formulated as calculating a pressure correction field and applying it to the physical hydrostatic reconstruction:
\begin{equation}
P_\text{corr} = P_\text{obs} + f_\theta(\mathbf{x}, \mathbf{p}, t) \ ,
\label{eq:pressure_correction}
\end{equation}
where $f_\theta$ represents the proposed PINF parameterized by weights $\theta$. The Neural Field predicts the continuous pressure residual $\delta P$ using spatiotemporal coordinates $\mathbf{x} = (\phi, \lambda, z)$, temporal mapping $t$, and physical features $\mathbf{p}$ (importantly including the sensor pressure bias). The final predicted height $h_\text{pred}$ is strictly routed through the un-deformed Hypsometric Equation using the corrected pressure $P_\text{corr}$:
\begin{equation}
h_\text{pred} = H_\text{ref} + \frac{R_\text{dry} T_\text{v}}{g_0} \ln\left(\frac{P_\text{ref}}{P_\text{corr}}\right) \ .
\label{eq:problem_formulation}
\end{equation}

\begin{figure}[htbp]
    \centering
    % Ensure your assembled image is named fig3_framework.png and placed in the image folder
    \includegraphics[width=1\linewidth]{image/fig3.pdf}
    \caption{The overall algorithmic framework. The Physical Baseline Branch derives a deterministic geometric height ($\hat{h}_{phy}$) based on the hypsometric equation and geoid undulation. Concurrently, the Neural Residual Branch leverages multi-resolution hash encoding and engineered terrain features to learn the complex spatial altitude residual ($\Delta h$). The final prediction ($h_{pred}$) is supervised by GNSS ground truth, with backpropagation (dashed line) restricted to the neural weights to preserve physical interpretability.}
    \label{fig:framework}
\end{figure}

\subsection{Spatiotemporal Encoding and Physics Bias Formulation}
\label{subsec:hash_and_terrain}

To effectively learn the highly non-linear residual field $\delta P$, the neural network must be capable of resolving high-frequency spatial variations induced by urban microclimates without heavily overfitting to single sensors. Standard coordinate-based Multi-Layer Perceptrons (MLPs) inherently suffer from ``spectral bias'', favoring low-frequency functions and struggling to fit sharp spatial gradients. To overcome this, we adapt the multi-resolution hash encoding technique~\cite{muller2022instant} alongside a novel temporal Fourier encoding map.

\subsubsection{Multi-Resolution Spatiotemporal Encoding}
Unlike traditional positional encoding that relies on fixed sinusoidal functions, hash encoding maps the spatial coordinates $(x, y, z)$ to trainable feature vectors across multiple resolution scales. For a given resolution level $l \in \{1, 2, \dots, L\}$ with a corresponding grid resolution $N_l$, we define a hash table $\mathcal{H}_l$ that maps spatial indices to continuous feature vectors:
\begin{equation}
\mathbf{e}_l(x, y, z) = \mathcal{H}_l\left(\text{hash}\left(\lfloor x \cdot N_l \rfloor, \lfloor y \cdot N_l \rfloor, \lfloor z \cdot N_l \rfloor\right)\right)
\label{eq:hash_lookup}
\end{equation}
where $\mathbf{e}_l \in \mathbb{R}^F$ is the feature vector with dimension $F=4$, and we utilize $L=16$ independent levels.

The grid resolutions follow a geometric progression to ensure a balanced receptive field:
\begin{equation}
N_l = N_{\text{min}} \cdot b^l, \quad b = \left(\frac{N_{\text{max}}}{N_{\text{min}}}\right)^{\frac{1}{L-1}}
\label{eq:resolution}
\end{equation}
with $N_{\text{min}} = 16$ and $N_{\text{max}} = 512$. The final spatial encoding $\phi(x, y, z)$ is the concatenation of features interpolated from all $L$ levels:
\begin{equation}
\phi(x, y, z) = [\mathbf{e}_1, \mathbf{e}_2, \dots, \mathbf{e}_L] \in \mathbb{R}^{L \times F}
\label{eq:concat_encoding}
\end{equation}
Furthermore, to model periodic atmospheric variations, the temporal observation variable $t$ is projected into a high-dimensional space via Fourier transformations across 6 distinct frequencies.

\subsubsection{Decoupling Hardware via Physics-Derived Pressure Bias}
The single greatest challenge for deep-learning-based metrology is generalizing to completely unseen uncalibrated sensor nodes. Direct MLPs will inherently memorize the systematic hardware bias of training units. To enforce zero-shot generalization and abstract the network away from strict node identities, we propose a Physics-Derived Pressure Bias formulation point:
\begin{equation}
P_\text{bias} = P_\text{obs} - P_\text{expected}
\label{eq:bias_formulation}
\end{equation}
where $P_\text{expected}$ is the ideal barometric reading analytically derived by inversing the uncalibrated GNSS altitude ($h_\text{GNSS}$) against the deterministic Hypsometric physical baseline.

By calculating $P_\text{bias}$ and feeding it as a learned embedding rather than using discrete sensor-ID tokens or vulnerable local terrain densities, the network fundamentally acts on the relative discrepancy space. To predict the pressure correction $\delta P$, this explicit physics characteristic is concatenated with the spatial hash encoding, temporal Fourier encoding, temperature $T$, and relative humidity $RH$ to form the complete continuous feature vector:
\begin{equation}
\mathbf{f}_{\text{input}} = [\phi(x, y, z), \phi_{\text{Fourier}}(t), P_\text{bias}, T, RH]
\label{eq:full_input}
\end{equation}
This fused vector $\mathbf{f}_{\text{input}}$ is subsequently processed by the Neural Residual Field's MLP to output the highly precise pressure correction $\delta P$.

\subsection{Altitude-Based Curriculum Learning Strategy}
\label{subsec:curriculum_learning}

Training continuous Neural Fields on massive and spatially heterogeneous IoT data presents severe convergence challenges. In an urban VLL environment, measurements acquired near ground-level arrays exhibit strong spatial correlation and are less perturbed by altitude-dependent atmospheric degradation, making them mathematically ``easier'' to fit. Conversely, measurements extending heavily vertically into the atmosphere are sparse, boundary-defining constraints prone to non-linear physical drift, representing ``harder'' samples. To algorithmically prevent the network from overfitting to densely sampled ground-level data while failing at higher altitudes, we propose a strict Altitude-Based Curriculum Learning structure.

\subsubsection{Three-Stage Curriculum Scheme}
Inspired by cognitive learning principles \cite{bengio2009curriculum}, the training process is explicitly structured to expose the SIREN-based geometry network to progressively more challenging geospatial distributions. We dynamically partition the full training dataset $\mathcal{D}$ into three distinct curriculum stages bound by absolute geometric altitude ($h$):
\begin{itemize}
    \item \textbf{Stage 1 (Easy - Foundation Formulation)}: The network is exclusively trained on low-altitude samples yielding stable thermodynamic behaviors. $\mathcal{D}_1 = \{ \mathbf{x} \in \mathcal{D} \mid h < 100\text{m} \}$. This safely establishes fundamental spatial hash structures.
    \item \textbf{Stage 2 (Medium - Moderate Expansion)}: The distribution bounds are relaxed to introduce mild-altitude geometries and early extrapolation behaviors. $\mathcal{D}_2 = \{ \mathbf{x} \in \mathcal{D} \mid h < 200\text{m} \}$.
    \item \textbf{Stage 3 (Hard - Extreme Outliers)}: The network is finally optimized against the entire dataset, introducing severe, high-altitude outliers. $\mathcal{D}_3 = \mathcal{D}$.
\end{itemize}

During training, the model systematically transitions from $\mathcal{D}_1 \rightarrow \mathcal{D}_2 \rightarrow \mathcal{D}_3$. Stage 1 and 2 operate for 30 epochs each, constructing stable global gradients, before Stage 3 fine-tunes on the complete unconstrained atmospheric volume for 80 epochs. If directly exposed to extreme vertical disparities without foundation (No Curriculum Strategy), gradient explosion fundamentally limits valid generalization to 9.27m MAE. By employing the 3-stage altitude scheme, the optimized framework confidently stabilizes across all bounds.

\begin{figure}
    \centering
    \includegraphics[width=1\linewidth]{image/fig5_curriculum.png}
    \caption{Curriculum learning strategy effectively preventing gradient collapse and stabilizing spatial hash geometries by sequencing altitude complexity.}
    \label{fig:curriculum}
\end{figure}

\subsubsection{SIREN Architecture and Optimization}
The fused feature vector $\mathbf{f}_{\text{input}}$ is linearly projected through a 3-layer explicit Multi-Layer Perceptron (hidden dimension=256) to output the localized $\delta P$. Crucially, instead of standard ReLU/SiLU components, we utilize Sinusoidal Representation Networks (SIREN) for the activation layers. SIREN architectures inherently preserve analytical derivative stability for periodic physical signals, preventing vanishing gradient stagnation over harsh geographical boundaries.

The forward pass is evaluated continuously per batch using an AdamW optimizer iteratively minimizing the exact residual between reconstructed orthometric height and strict observed GNSS MSL:
\begin{equation}
\mathcal{L}(\theta) = \frac{1}{B} \sum_{i=1}^B \left| h^\text{pred}_{i} - h^{true}_i \right|
\label{eq:loss_function}
\end{equation}
Furthermore, a Cosine Annealing learning rate scheduler with warm restarts is initialized ($10^{-3}$), effectively allowing periodic deep landscape exploration and preventing premature optimization stagnation.